\begin{table}[H]
\centering
\caption{Approximate Peer/Rank Diagnostic by Score Density}
\label{tab:density_decomp}
\begin{threeparttable}
\begin{tabular}[t]{lccc}
\toprule
 & \multicolumn{3}{c}{Score-density tercile} \\
\cmidrule(l{3pt}r{3pt}){2-4}
 & (1) Low & (2) Medium & (3) High \\
 & (tails) & & (middle) \\
\midrule
\multicolumn{4}{l}{\textit{Panel A: Approximate peer/rank estimate by density tercile}} \\
Peer mean EB ability ($\hat{\zeta}$) &  -0.179** &  -0.099* &  -0.123** \\
 & (  0.081) & (  0.053) & (  0.062) \\
N & 1708 & 1672 & 1574 \\
\addlinespace
\midrule
 & \multicolumn{3}{c}{Pooled interaction} \\
\cmidrule(l{3pt}r{3pt}){2-4}
\multicolumn{4}{l}{\textit{Panel B: Interaction specification}} \\
Peer mean EB ability ($\hat{\zeta}$) & \multicolumn{3}{c}{ -0.141} \\
Peer EB $\times$ density interaction & \multicolumn{3}{c}{ -0.024} \\
 & \multicolumn{3}{c}{(  0.039)} \\
\bottomrule
\end{tabular}
\begin{tablenotes}[para,flushleft]
\scriptsize
\setlength{\emergencystretch}{2em}
\item \textit{Notes:} Panel~A reports the approximate Borusyak--Hull control function peer/rank coefficient ($\hat{\zeta}$) separately for terciles of the baseline score density evaluated at each student's own score. Students in the tails of the distribution (low density) have few classmates at similar ability levels, so their within class rank is insensitive to changes in peer composition. Students in the middle (high density) have many nearby classmates, so rank shifts more with peer changes. In the stylized rank model, $\hat{\zeta} \approx \pi + \eta \cdot (-f(\theta_i))$, where $\pi$ is the mean peer effect, $\eta$ is the rank effect, and $\,f(\theta_i)$ is the score density. If rank effects matter in this reduced-form composite ($\eta \neq 0$), $\hat{\zeta}$ would tend to be more negative in the high-density tercile, where rank is more responsive to peer shifts. Panel~B reports a pooled interaction of peer mean with standardized density; the interaction is a diagnostic for whether the peer/rank coefficient changes with local score density. All specifications include baseline-decile $\times$ treatment $\times$ grade FE and strata FE, so they should be read as approximate cell-control specifications rather than exact design-based recentering checks. Density is estimated using a Gaussian kernel with Silverman bandwidth, computed within grade, and standardized to mean zero and unit variance within grade. Standard errors are clustered at the school level. ***, **, and * indicate significance at 1\%, 5\%, and 10\%.
\end{tablenotes}
\end{threeparttable}
\end{table}
